%&latex
\documentclass[a4paper]{article}

\frenchspacing

\usepackage[cp1250]{inputenc}
\usepackage[czech]{babel}

\usepackage{a4wide}
\usepackage{amsmath, amsthm, amssymb, amsfonts}
\usepackage[mathcal]{eucal}
\usepackage{graphicx}
\usepackage{url}
\usepackage{color}
\usepackage{wrapfig}
\usepackage{capt-of}
\usepackage{float}
\usepackage{ifthen}
\usepackage{subfig}
\usepackage[normalem]{ulem}

% sirka a vyska textu nastavena jako papir, vsechny okraje vynulovany a pridano 20pt na kazdou stranu
% horizontalni rozmery
\setlength{\textwidth}{\paperwidth}
\addtolength{\textwidth}{-40pt}
\addtolength{\hoffset}{-1in}
\addtolength{\hoffset}{20pt}
\setlength{\oddsidemargin}{0in}
\setlength{\marginparsep}{0in}
% vertikalni rozmery
\setlength{\textheight}{\paperheight}
\addtolength{\textheight}{-60pt}
\addtolength{\voffset}{-1in}
\addtolength{\voffset}{20pt}
\setlength{\topmargin}{0in}
\setlength{\headheight}{0in}
\setlength{\headsep}{0in}


%Obrazek na miste
%pouziti
%%\obrazeknahore{adresa}{popisek}{label}
\long\def\obrazeknahore#1#2#3 {

\begin{figure}[t]
    \centering
    \includegraphics[width=0.8\textwidth]{#1}
    
    \caption{#2}
    \label{#3}
    
\end{figure}

}


%==========================================
%PEKELNA MAKRA NA ZAROVNANI OBRAZKU DOPRAVA

\makeatletter


%tohle je makro, ktere mi dovoluje obtekani i u kratkych environmentu
%ABSOLUTNE nechapu, jak to funguje, ale funguje to
%viz http://tex.stackexchange.com/questions/26078/ 
\def\odrovnej{\@@par
\ifnum\@@parshape=\z@ \let\WF@pspars\@empty \fi % reset `parshape'
\global\advance\c@WF@wrappedlines-\prevgraf \prevgraf\z@
\ifnum\c@WF@wrappedlines<\tw@ \WF@finale \fi}

\makeatother



%---
%makro, co da obrazek doprava a ostatni text ho obteka
%(bez toho predchazejiciho makra to ale poradne nebeha)
%pouziti:
%\obrazekvpravo{adresa}{popisek}{label}{procento sirky}
\long\def\obrazekvpravo#1#2#3#4{

\setlength\intextsep{-20pt}

    \begin{wrapfigure}{r}{#4\textwidth}
      \begin{center}
          \vspace{-10pt}
          
        \includegraphics[width=#4\textwidth]{#1}
        \vspace{-10pt}
        
      \end{center}
      
      \caption{#2}
      \label{#3}
      
      
    \end{wrapfigure}

\setlength\intextsep{0pt}

    
}


\long\def\obrazekvpravonocap#1#2#3#4#5{
  \ifthenelse{\equal{#4}{0}}{\setlength\intextsep{-5pt}}{\setlength\intextsep{#4pt}}
      \begin{wrapfigure}{r}{#3\textwidth}
        \begin{center}
          \vspace{-10pt}            
          \includegraphics[width=#3\textwidth]{#1}
          \vspace{-10pt}
        \end{center}
        \label{#2}
      \end{wrapfigure}
  \ifthenelse{\equal{#5}{0}}{\setlength\intextsep{0pt}}{\setlength\intextsep{#5pt}}
}

%---
%makro pro pripady, kdy wrapfigure neco mrsi
%je to docela pekelne
%je nutne mu dat jak text vpravo, tak text vlevo
%a nevim, jestli bude 100% fungovat, ale doufam, ze jo

%pouziti:
%\obrazekvpravominipage{adresa}{popisek}{label}{procento sirky}{1 - procento sirky}{text vlevo}
\long\def\obrazekvpravominipage#1#2#3#4#5#6{

\noindent\begin{minipage}{#5\linewidth}
\vspace{0pt}
#6
\end{minipage}
\hspace{0.5cm}
\noindent\begin{minipage}{#4\linewidth}
\vspace{0pt}
\centering
\includegraphics[width=0.9\textwidth]{#1}
\captionof{figure}{#2}
\label{#3}
\end{minipage}

}

%KONEC PEKELNYCH MAKER
%=====================

\def\lebka{\includegraphics[width=3mm]{../common/lebka}}

% makra pro poznamku u vyrokove a predikatove logiky
\def\vl{ -- ve v�rokov� logice}
\def\pl{ -- v predik�tov� logice}


%Vacsina prostredi je dvojjazicne. V pripade, ze znenie napr pozorovania je pisane po slovensky, malo by byt po slovensky aj oznacenie.

\def\ifis#1#2{\ifthenelse{\equal{#1}{0}}{}{#2}}

%\newenvironment{e}[3]{\pagebreak[2]\noindent\textbf{\ifis{#2}{$\bigstar$}#1}\ifis{#3}{\emph{~(#3)}}\par\noindent\leftskip 10pt }{\odrovnej\par\bigskip}
\newenvironment{e}[3]{\pagebreak[2]\noindent\textbf{\ifis{#2}{\lebka~}#1}\ifis{#3}{\emph{~(#3)}}\ifis{#2}{~\lebka}\par\noindent\leftskip 10pt }{\odrovnej\par\bigskip}


%obecne prostredie, ktore ma vyuzitie pri specialnych odstavcoch ako (uloha, algoritmus...) aby nevzniklo dalsich x prostredi
\newenvironment{obecne}[1]{\pagebreak[2]\noindent\textbf{#1}\par\noindent\leftskip 10pt}{\odrovnej\par\bigskip}

\newenvironment{report}{\pagebreak[2]\noindent\textbf{Report}\em\par\noindent\leftskip 10pt}{\par\bigskip}

%\newenvironment{reportN}[1]{\pagebreak[2]\noindent\textbf{Report~}\emph{(#1)}\emph\par\noindent\leftskip 10pt}{\odrovnej\par\bigskip}
\newenvironment{reportN}[1]{\pagebreak[2]\noindent\textbf{Report~}\emph{(#1)}\em\par\noindent\leftskip 10pt}{\odrovnej\par\bigskip}

\newenvironment{penumerate}{
\begin{enumerate}
  \setlength{\itemsep}{1pt}
  \setlength{\parskip}{0pt}
  \setlength{\parsep}{0pt}
  %\setlength{\topsep}{200pt}
  \setlength{\partopsep}{200pt}
}{\end{enumerate}}

\def\pismenka{\numberedlistdepth=2} %pouzit, ked clovek chce opismenkovany zoznam...

\newenvironment{pitemize}{
\begin{itemize}
  \setlength{\itemsep}{1pt}
  \setlength{\parskip}{0pt}
  \setlength{\parsep}{0pt}
}{\end{itemize}}

\definecolor{gris}{gray}{0.95}
\newcommand{\ramcek}[2]{\begin{center}\fcolorbox{white}{gris}{\parbox{#1}{#2}}\end{center}\par}


\title{\LARGE U�ebn� texty k st�tn� bakal��sk� zkou�ce \\ Matematika \\ Z�klady teorie funkc� v�ce prom�nn�ch}

\begin{document}


\setcounter{section}{4}
%&latex
\section{Z�klady teorie funkc� v�ce prom�nn�ch}

\begin{e}{Po�adavky}{0}{0}
\begin{pitemize}
        \item Parci�ln� derivace a tot�ln� diferenci�l
        \item V�ty o st�edn� hodnot�
        \item Extr�my funkc� v�ce prom�nn�ch
        \item V�ta o implicitn�ch funkc�ch
\end{pitemize}
\end{e}

\subsection{Parci�ln� derivace a tot�ln� diferenci�l}
\def\bodx{\boldsymbol{\mathrm{x}}}
P�i definici derivace pro funkce v�ce prom�nn�ch m��eme postupovat n�kolika zp�soby. Prvn� mo�nost je uva�ovat hodnoty dan� funkce pouze na jist� p��mce, konkr�tn� na p��mk�ch rovnob�n�ch se sou�adn�mi osami. Z�sk�me tak v podstat� funkci jedn� prom�nn�, kterou ji� derivovat um�me. Nev�hodou v�ak je, �e tato derivace (naz�van� parci�ln�) poskytuje informaci o chov�n� funkce pouze v��i velice speci�ln�m podmno�in�m defini�n�ho oboru (p��mk�ch rovnob�n�ch se sou�adn�mi osami).
\\
\obrazekvpravonocap{matematika/obrazky/Parcialni_derivace_graf_vyznam}{fig:parc_derivace}{0.35}{44}{-50}
\begin{e}{Definice}{1}{Parci�ln� derivace}
Nech� $f: \mathbb{R}^n \rightarrow \mathbb{R}$, $\bodx
= [x_1, \ldots, x_n] \in \mathbb{R}^n$. Potom \textit{parci�ln� derivac� funkce $f$ podle i-t� slo�ky v bod� \boldsymbol{x}} nazveme limitu
$$\frac{\partial f}{\partial x_i}(\bodx) = \displaystyle\lim_{h \rightarrow 0}~\frac{f(x_1,\ldots,x_{i-1},x_i + h,x_{i+1},\ldots,x_n) - f(x_1,\ldots,x_n)}{h}$$
pokud tato limita existuje. Parci�ln� derivace m��e b�t vlastn� i nevlastn�.
\end{e}
\begin{e}{Pozn�mka}{0}{Geometrick� v�znam parci�ln� derivace}
Stejn� jako u funkc� jedn� prom�nn� p�edstavuje parci�ln� derivace funkce $f$ podle j-t� prom�nn� v bod� $\bodx$ sm�rnici te�ny ve sm�ru osy $x_j$.
\\\\
Pro jednoduchost uva�ujme 2D probl�m. Parci�ln� derivace funkce $f : \mathbb{R}^{2}\rightarrow
\mathbb{R}$ v bod� $[x_0; y_0]\in \mathbb{R}^{2}$, v n�m� m� funkce hodnotu $z_0 = f(x_0,y_0)$, odpov�daj� sm�rnic�m\footnote{p�ipome�me �e sm�rnice
je definovan� jako: $m = \frac{\Delta y}{\Delta x} = \frac{\text{rise}}{\text{run}}
=tg\alpha$ kde $\alpha$ je �hel sv�ran� s kladnou ��st� osy $x$. } te�en ve sm�ru jednotliv�ch sou�adnicov�ch os. Sm�rnici te�ny $t_1$ tedy z�sk�me jako parci�ln� derivaci podle x, tzn. $\operatorname{tg} \alpha = \frac{\partial f}{\partial x}(x_0,y_0)$, a podobn� pro sm�rnici te�ny $t_2$ dostaneme $\operatorname{tg} \beta = \frac{\partial f}{\partial y}(x_0,y_0)$.
\\Parci�ln� derivace funkce $f : \mathbb{R}^{n}\rightarrow
\mathbb{R}$ v bod� $\bodx\in \mathbb{R}^{n}$ jsou derivace ve sm�rech
rovnob�n�ch s n�kterou sou�adnicovou osou. Jejich zobecn�n�m jsou derivace ve sm�ru.
\end{e}
\begin{e}{V�ta\footnote{Tyto vzorce fakticky jsou vzorce pro funkce jedn� promenn�, proto�e $\partial$ se poc�t� z funkce z�visej�c� jen na $x_i$,
p�i�em� ostatn� prom�nn� jsou br�ny jako konstanty (a derivace konstanty je nula). Podobn� pro diferenci�ly.}}{0}{Aritmetika parci�ln� derivace}
Nech� $c,d\in\mathbb{R}$ a $f, g: \mathbb{R}^n \rightarrow \mathbb{R}$ maj� v bod� $\bodx \in \mathbb{R}^n$  vlastn� derivaci podle $x_i$. Potom:
\begin{pitemize}
\item $\frac{\partial (cf\pm dg)}{\partial x_i}(\bodx) =c\frac{\partial f}{\partial x_i}(\bodx) \pm d\frac{\partial g}{\partial x_i}(\bodx)$

\item $\frac{\partial (fg)}{\partial x_i}(\bodx) =\frac{\partial f}{\partial x_i}(\bodx)g(\bodx) + f(\bodx)\frac{\partial g}{\partial x_i}(\bodx)$


\item $\frac{\partial (f\div g)}{\partial x_i}(\bodx) = \frac{\frac{\partial f}{\partial x_i}(\bodx)g(\bodx) - f(\bodx)\frac{\partial g}{\partial x_i}(\bodx)}{g^2(\bodx)}$,
pokud $g(\bodx)\neq 0$
\end{pitemize}
\end{e}

\begin{e}{Definice}{1}{Derivace ve sm�ru vektoru}
Nech� $f: \mathbb{R}^n \rightarrow \mathbb{R}$, $v \in \mathbb{R}^n\setminus\{0^n\}$
je nenulov� vektor, $\bodx = [x_1, \ldots, x_n] \in \mathbb{R}^n$. Potom \textit{derivac� funkce $f$ ve sm�ru vektoru v} nazveme limitu
$$\frac{\partial f}{\partial v}(\bodx) = \displaystyle\lim_{h \rightarrow 0}~\frac{f(\bodx + h \cdot v) - f(\bodx)}{h}$$
pokud tato limita existuje.
\end{e}

\begin{e}{Pozn�mka}{0}{Geometrick� v�znam sm�rov� derivace}
Pokud je derivace ve sm�ru kladn�, funkce roste v tomto sm�ru, opa�n� je-li derivace ve sm�ru z�porn�, funkce kles� v tomto sm�ru. Je-li sm�rov� vektor nav�c jednotkov� ud�v� n�m derivace ve sm�ru rychlost r�stu/kles�n� v dan�m sm�ru.
\\\\
P�edstavte si nap��klad oblast v t��rozm�rn�m euklidovsk�m prostoru, kde funkce $f$ m��� teplotu a kterou prol�t� po p��mo�ar� dr�ze ��stice.
Sm�rov� derivace pak ud�v� okam�itou zm�nu teploty ��stice ve chv�li,
kdy se nach�z� v bodu $\bodx$ a m� vektor rychlosti $v$.
\end{e}

\obrazekvpravonocap{matematika/obrazky/gradient}{fig:parc_derivace}{0.2}{0}{0}
\begin{e}{Definice}{1}{Gradient}
Nech� $f: \mathbb{R}^n \rightarrow \mathbb{R}$, $\bodx = [x_1, \ldots, x_n] \in \mathbb{R}^n$ a nech� existuj� v�echny parci�ln� derivace funkce $f$ v bod� $\bodx$ a jsou vlastn�. Pak vektor $$\nabla f(\bodx) = [\frac{\partial f}{\partial x_1}(\bodx),\ldots,\frac{\partial f}{\partial x_n}(\bodx)]$$ naz�v�me \textit{gradientem funkce $f$ v bod� $\bodx$}.
\end{e}

\begin{e}{Pozn�mka}{0}{Geometrick� v�znam gradientu}
Uva�ujme modelovou �lohu ve 2D, kde funkce $f (x; y)$ popisuje reli�f krajiny
(3D funkce).
Gradient v bod� x lze interpretovat jako vektor kolm� na vrstevnici grafu funkce v dan�m bod� x. Z fyzik�ln�ho hlediska gradient funkce ur�uje sm�r nejv�t��ho stoup�n�.
\end{e}

\begin{e}{Definice}{1}{Tot�ln� diferenci�l}
Nech� $f: \mathbb{R}^n \rightarrow \mathbb{R}$, $\bodx = [x_1, \ldots, x_n] \in \mathbb{R}^n$ a nech� $v \in \mathbb{R}^n$. Existuje-li line�rn� zobrazen� $Df(\bodx)(v)$ takov�, �e plat�:
$$\displaystyle\lim_{\|v\| \rightarrow 0}~\frac{f(\bodx + h) - f(\bodx) - Df(\bodx)(v)}{\|v\|} = 0$$
\noindent potom toto zobrazen� naz�v�me \textit{tot�ln� diferenci�l funkce $f$ v bod� $\bodx$}.
\end{e}

\begin{e}{Pozn�mka}{0}{Geometrick� v�znam tot�ln�ho diferenci�lu}
Tot�ln� diferenci�l $Df(\bodx)(v)$ n�m poskytuje line�rn� aproximaci obecn� funkce f (\bodx) v okol� dan�ho bodu, nap�. ve 2D obecn�
neline�rn� funkci f(\bodx) nahrazujeme line�rn� fukc� (rovinou).
\end{e}

\begin{e}{Definice}{0}{Parci�ln� derivace druh�ho ��du}
Nech� $M \subseteq \mathbb{R}^n$ otev�en�, a nech� m� funkce $f$ parci�ln� derivaci $\frac{\partial f}{\partial x_i}$. Pak pro $\bodx \in M$ definujeme \textit{parci�ln� derivaci druh�ho ��du (podle i-t� a j-t� slo�ky)} jako $\frac{\partial^2f}{\partial x_i \partial x_j}(\bodx) = \frac{\partial}{\partial x_j} (\frac{\partial f}{\partial x_i}(\bodx))$.
\end{e}

\begin{e}{Definice}{0}{Druh� diferenci�l}
Nech� $f: \mathbb{R}^n \rightarrow \mathbb{R}$ a $\bodx \in \mathbb{R}^n$. �ekneme, �e $f$ m� v bod� $\bodx$ \textit{druh� diferenci�l}, pokud ka�d� parci�ln� derivace $f$ m� v bod� $\bodx$ tot�ln� diferenci�l. Druh� diferenci�l je biline�rn� zobrazen� $D^2f(\bodx): \mathbb{R}^n \times \mathbb{R}^n \rightarrow \mathbb{R}$ a m� tedy n�sleduj�c� tvar:
$$D^2f(\bodx)(h, k) = \sum_{i=1}^n\sum_{j=1}^n\frac{\partial^2f}{\partial x_i \partial x_j}(\bodx)h_ik_j$$
Pou�ijeme-li analogii gradientu pro prvn� diferenci�l, m��eme ��ct, �e druh� diferenci�l je reprezentov�n matic�:
\begin{equation}
\left(
\frac{\partial^2 f}{\partial x_i \partial x_j}(\bodx)
\right)
^{n, n}
_{i = 1, j = 1}
\end{equation}
\end{e}

\begin{e}{V�ta}{1}{Diferenci�l implikuje parci�ln�/sm�rov� derivace a spojitost  
\scriptsize(Klazar 2.1-3)\normalsize}
Nech� $f: \mathbb{R}^n \rightarrow \mathbb{R}$ m� v bod� $\bodx \in \mathbb{R}^n$ tot�ln� diferenci�l. Potom:
\begin{pitemize}
\item pro $\forall v \in \mathbb{R}^n\setminus\{0^n\}$ existuje $D_vf(\bodx)$ vlastn� a plat� $D_vf(\bodx) = Df(\bodx)(v)$.
\item existuj� v�echny parci�ln� derivace a pro $\forall v \in \mathbb{R}^n: Df(\bodx)(v) = \sum_{i=0}^n\frac{\partial f}{\partial x_i}(\bodx) \cdot v_i= \left<\nabla f(\bodx), v\right>$).
\item $f$ je spojit� v $\bodx$.
\end{pitemize}
\end{e}

\begin{e}{V�ta}{1}{Spojit� parci�ln� derivace implikuj� diferenci�l \scriptsize(Klazar 2.4)\normalsize}
Nech� $f: \mathbb{R}^n \rightarrow \mathbb{R}$ m� v bod� $\bodx \in \mathbb{R}^n$ spojit� v�echny parci�ln� derivace. Potom m� $f$ v bod� $\bodx$ tot�ln� diferenci�l.
\end{e}

\begin{e}{V�ta}{0}{Posta�uj�c� podm�nka pro existenci druh�ho diferenci�lu}
Nech� $M \subseteq \mathbb{R}^n$ je otev�en� a $f$ m� spojit� parci�ln� derivace druh�ho ��du na M. Potom $f$ m� v ka�d�m bod� z $M$ druh� diferenci�l.
\end{e}

\begin{e}{V�ta}{0}{Aritmetika tot�ln�ho diferenci�lu \scriptsize(Klazar 2.7)\normalsize}
Nech� $c,d\in\mathbb{R}$ a $f, g: \mathbb{R}^n \rightarrow \mathbb{R}$ maj� v bod� $\bodx \in \mathbb{R}^n$ tot�ln� diferenci�l.  Potom existuj� tot�ln� diferenci�ly $ D(cf \pm dg)(\bodx)$, $D(f \cdot g)(\bodx)$. Pokud nav�c $g(\bodx) \neq 0$ existuje i $D(f \div g)(\bodx)$. Nav�c plat�:
\begin{pitemize}
\item $D(cf \pm dg)(\bodx) = cDf(\bodx) \pm dDg(\bodx)$, pokud $\exists D(cf \pm dg)(\bodx)$
\item $D(f \cdot g)(\bodx) = Df(\bodx)g(\bodx) + f(\bodx)Dg(\bodx)$.
\item $D(f \div g)(\bodx) = \frac{Df(\bodx)g(\bodx) - f(\bodx)Dg(\bodx)}{g^2(\bodx)}$.
\end{pitemize}
\end{e}

\begin{e}{V�ta}{0}{Diferenci�l slo�en�ho zobrazen�  \scriptsize(Klazar 2.8
+ d�sledky)\normalsize}
M�jme funkci $f: \mathbb{R}^n \rightarrow \mathbb{R}$ a $n$ funkc� $g_j: \mathbb{R}^m \rightarrow \mathbb{R}$. Nech� $a \in \mathbb{R}^m$ a $b \in \mathbb{R}^n$ a $b_j = g_j(a)$. Nech� existuj� $Df(a)$ a $Dg_i(a), i = 1 \ldots n$. Definujeme-li zobrazen� $H: \mathbb{R}^m \rightarrow \mathbb{R}$ p�edpisem $H(x) = f(g_1(x), \ldots, g_n(x))$, potom $H$ m� v bod� $a$ tot�ln� diferenci�l a pro $h \in \mathbb{R}^m$ plat�
$$DH(a)(h) = \sum_{i = 1}^n\left(\sum_{j = 1}^n\frac{\partial f}{\partial y_j}(b)\frac{\partial g_j}{\partial x_i}(a)\right)h_i$$
Z �eho� plyne tzv. \emph{�et�zkov� pravidlo}, tj.:
$$\frac{\partial H}{\partial x_i}(a)=\sum_{j=1}^n \frac{\partial f}{\partial y_j}(b)\frac{\partial g_j}{\partial x_i}(a)$$
\end{e}

\begin{e}{V�ta}{0}{Z�m�nnost parci�ln�ch derivac� druh�ho ��du  \scriptsize(Klazar 2.9)\normalsize}
M�jme funkci $f: \mathbb{R}^n \rightarrow \mathbb{R}$. Nech� $f$ m� spojitou parci�ln� derivaci $\frac{\partial^2f}{\partial x_i \partial x_j}(\bodx)$. Potom existuje i $\frac{\partial^2f}{\partial x_j \partial x_i}(\bodx)$ a ob� tyto parci�ln� derivace druh�ho ��du se rovnaj�.
\end{e}

\begin{e}{D�sledek}{0}{0}
D�sledkem dvou pr�v� uveden�ch v�t je fakt, �e matice, kter� reprezentuje druh� diferenci�l funkce $f$ v bod� $\bodx$ (tedy hovo��me o situaci, kdy $f$ m� v bod� $\bodx$ druh� diferenci�l), je symetrick�.
\end{e}
\begin{e}{Report}{0}{Fiala} Parci�ln� derivace a tot. diferenci�l. - Uvedl jsem definice, jak to souvis� se spojitost� funkce a jak se navz�jem implikuj�, geometrick� v�znam a nutnou podm�nku 1. ��du pro extr�m. V�e bez d�kazu. Po tro�e dokop�v�n� ze m� dostal, �e gradient je kolm� na vrstevnici v dan�m bod� a kdy� se m� ptal, jestli bych to um�l dok�zat a j� �ekl, �e ne, odpov�d�l: "No vid�te, to jste se taky v anal�ze u�ili :-)" nakonec za 1-
\end{e}

\subsection{V�ty o st�edn� hodnot�}

\begin{e}{V�ta}{1}{Lagrangeova o st�edn� hodnot� pro funkce v�ce prom�nn�ch  \scriptsize(Klazar 2.5)\normalsize}
Nech� $f: \mathbb{R}^n \rightarrow \mathbb{R}$ a $a, b \in \mathbb{R}^n$. Nech� $f$   m� diferenci�l a je spojit� v ka�d�m vnit�n�m bod� �se�ky $(a, b)$. Potom $\exists$ bod $\xi \in (a, b)$ takov�, �e
$$f(b) - f(a) = D f(\xi)(b - a).$$

\begin{e}{D�kaz}{0}{0}
Plyne z Lagrangeovy v�ty o st�edn� hodnot� pro funkci $F:[0,1]\to\mathbb{R}$ definovanou p�edpisem $F(t)=f(a+t(b-a))$ a �et�zkov�ho pravidla.
\end{e}
\end{e}

\begin{e}{V�ta}{0}{Nulovost derivace implikuje konstatnost funkce \scriptsize(Klazar 2.6)\normalsize}
M�-li re�ln� funkce $m$ promenn�ch v ka�d�m bod� otev�en� a souvisl�\footnote{�ekneme, �e otevren� mno�ina $D$ v $\mathbb{R}^n$ je souvisl�, kdy� lze ka�d� jej� dva body spojit lomenou �arou, kter� cel� le�� v $D$.} mno�iny nulov� diferenci�l, je na t�to mno�ine konstantn�.
\end{e}

\begin{e}{D�sledek}{0}{Nulovost parci�ln� derivace implikuje konstatnost funkce \scriptsize(Klazar 2.6 D�sledek)\normalsize}
M�-li re�ln� funkce $m$ prom�nn�ch v ka�d�m bod� otev�en� a souvisl�
mno�iny ka�dou parci�ln� derivaci nulovou, je na t�to mno�ine konstantn�.
\end{e}

\subsection{V�ta o implicitn�ch funkc�ch}

\begin{e}{Definice}{0}{Implicitn� vs. explicitn� zad�n� funkce}
Nech� $f: \mathbb{R}^n \rightarrow \mathbb{R}$ je funkce zadan� ve tvaru:
$$y = f(x), x \in \mathbb{R}^n$$
Toto zad�n� (kdy je d�n vzorec vyjad�uj�c� hodnotu z�visle prom�nn� $y$ na hodnot� nez�visl� prom�nn� $x$), naz�v�me \textit{explicitn�m zad�n�m funkce}.
\\\\
N�kdy v�ak m��e b�t z�vislost $y$ na $x$ (tj. funkce $y$ prom�nn� $x$) vyj�d�ena pouze rovnic� s nulovou pravou stranou, ve kter� na lev� stran� vystupuj� jak $x$ (nez�visl� prom�nn�) tak $y$ (z�visl� prom�nn�)
$$F(x, y) = 0, x \in\mathbb{R}^n~ a~~ y\in\mathbb{R}$$
Toto zad�n� naz�v�me \textit{implicitn�m zad�n�m funkce}.
\\\\\footnotesize Na��m c�lem bude tedy nal�zt a popsat mno�inu
$M_I = \{f[x, y] \in \mathbb{R}^{n+1} : F(x, y) = 0\}$
a ur�it, kdy tato mno�ina $M_I$ popisuje n�jakou implicitn� funkci a za jak�ch
p�edpoklad� lze p�ej�t od implicitn� zadan� funkce k explicitn�mi popisu (tj. kdy lze jednu prom�nnou vyj�d�it jako funkci ostatn�ch prom�nn�ch).
\end{e}

\begin{e}{V�ta}{1}{O existenci implicitn� funkce (v $R^{n+1}$)  }
Nech� $F : \mathbb{R}^{n+1} \rightarrow \mathbb{R}$ je spojit� v okol� bodu $[\bodx_0, y_0]$, kde $\bodx_0 \in\mathbb{R}^{n}$ a $y_0 \in\mathbb{R}$. D�le
p�edpokl�dejme, �e funkce $F$ m� v okol� bodu $[\bodx_0, y_0]$ spojitou parci�ln� derivaci podle prom�nn� $y$ a plat�
\begin{pitemize}
\item $\frac{\partial F}{\partial y}(\bodx_0,y_0) \neq 0$,
\item $F(\bodx_0, y_0) = 0$.
\end{pitemize}
Potom rovnice $F(\bodx, y) = 0$ definuje na okol� bodu $[\bodx_0, y_0]$ implicitn� n�jakou funkc� $y = f(\bodx) = f (x_1, x_2, ...,x_n).$
\end{e}

\begin{e}{V�ta}{1}{O derivaci implicitn� funkce (v $R^{n+1}$)}
Nech� plat� p�edpoklady p�edchoz� v�ty a $F$ m� spojit� parci�ln� derivace a� do ��du $m$ v okol� bodu $[\bodx_0, y_0]$. Potom funkce $y = f (\bodx)$ (jednozna�n� ur�en� rovnic� $F(\bodx, y) = 0$) m� spojit� parci�ln� derivace a� do ��du $m$. Speci�ln� pro $m = 1$ plat�:
$$\frac{\partial f}{\partial x_i}(\bodx) = -\frac{\frac{\partial F}{\partial x_i}(\bodx, y)}{\frac{\partial F}{\partial y}(\bodx, y)}$$
\end{e}

\begin{e}{Pozn�mka}{0}{0}
Na tomto m�st� uvedeme malou, ale pro n�s d�le�itou pozn�mku z algebry. M�jme bod $a \in \mathbb{R}^n$ a funkce $F_j, j=1 \ldots n$, $F_j: \mathbb{R}^n \rightarrow \mathbb{R}$, kter� maj� v�echny sv� parci�ln� derivace. Potom determinant
\begin{equation}
JF_{j=1}^n(a)= \left|\frac{\partial (F_1, \ldots, F_n)}{\partial (x_1,\ldots,x_n)}\right| = \det
\left(
\frac{\partial F_i}{\partial x_j}(a)
\right)
_{i = 1, j = 1}^{n,n}
\end{equation}
nazveme Jakobi�nem funkc� $F_j$ (v bod� $a$) vzhledem k prom�nn�m $x_1, \ldots, x_n$. Pojem Jakobi�n lze ekvivalentn� zav�st i pomoc� vektorov�ch funkc�. To zde v�ak nebudeme pot�ebovat.
\end{e}

\begin{e}{V�ta}{0}{O implicitn�ch funkc�ch (p��pad v $R^{n + m}$) \scriptsize(Klazar 2.12)\normalsize}
Nech� $F_j: G \rightarrow \mathbb{R}, j=1 \ldots m$, kde $G \subseteq \mathbb{R}^{n+m}$ je otev�en� mno�ina. Uva�ujme body $x_0 \in \mathbb{R}^n$, $y_0 \in \mathbb{R}^m$ takov�, �e $[x_0, y_0] \in G$ a $F_j([x_0, y_0]) = 0$ pro v�echny $j=1 \ldots m$. Nech� ka�d� funkce $F_j$ m� spojit� parci�ln� derivace a nech� nav�c $JF_{j=1}^m([x_0, y_0]) \neq 0$. Potom existuje okol� $U \subseteq \mathbb{R}^n$ bodu $x_0$ a okol� $V \subseteq \mathbb{R}^m$ bodu $y_0$ takov�, �e pro $\forall x \in U$ existuje pr�v� jedno $y \in V$ takov�, �e $F_j([x,y]) = 0, j=1 \ldots m$. Nav�c, ozna��me-li $y_j = \varphi_j(x), j=1 \ldots m$, potom $\varphi_j$ m� spojit� parci�ln� derivace na $U$ a plat�:
$$\frac{\partial \varphi_i}{\partial x_j}(x) = -\frac{\left|\frac{\partial (F_1, \ldots, F_m)}{\partial (y_1,\ldots,y_{i-1},x_j,y_{i+1},\ldots,y_m)}\right|}{\left|\frac{\partial (F_1, \ldots, F_m)}{\partial (y_1,\ldots,y_m)}\right|}$$
\end{e}

\begin{e}{V�ta}{0}{O inverzn�ch funkc�ch}
D�sledkem v�ty o implicitn�ch funkc�ch je n�sleduj�c� v�ta:  Nech� $f : U \to \mathbb{R}^m$, kde $U \subseteq \mathbb{R}^m$ je okol� bodu $x_0$, je zobrazen� se spojit�mi parci�ln�mi derivacemi, kter� m� v $x_0$ nenulov� jakobi�n. Potom existuj� okol� $U_1 \subseteq U$ a $V \subseteq \mathbb{R}^m$ bod� $x_0$ a $y_0 = f(x_0)$ takov�, �e $f : U_1 \to V$ je bijekce, inverzn� zobrazen� $f^{-1} : V \to U_1$ m� spojit� parci�ln� derivace a pro ka�d� $x \in U_1$ v bod� $y = f(x) \in V$ m�me
$$Df^{-1}(y) = (Df(x))^{-1}$$
Jacobiho matice zobrazen� $f^{-1}$ v bod� $y$ je tedy inverzn� k Jacobiho matici zobrazen� $f$ v bod� $x$.
\end{e}
\begin{e}{Report}{0}{Zahradn�k} implicitni fce - o implicnich fcich sem si skoro neskrtl (jedna blba veta ve vic variantach a kdyz si clovek nevzpomene je v ... nakonec za 2
\end{e}
\subsection{Extr�my funkc� v�ce prom�nn�ch}

\begin{e}{Definice}{0}{Extr�my funkce}
Nech� $f: \mathbb{R}^n \rightarrow \mathbb{R}$, $\overline{X} \in \mathbb{R}^n$, $M \subseteq \mathbb{R}^n$. �ekneme, �e bod $\overline{X}$ je bodem \textit{maxima funkce $f$ na mno�in� $M$}, pokud $\forall X \in M: f(\overline{X}) \geq f(X) $. Analogicky definujeme \textit{minimum funkce $f$ na mno�in� $M$}.
\end{e}

\begin{e}{Definice}{0}{Lok�ln� extr�my funkce}
Nech� $f: \mathbb{R}^n \rightarrow \mathbb{R}$, $\overline{X} \in \mathbb{R}^n$, $M \subseteq \mathbb{R}^n$. �ekneme, �e bod $\overline{X}$ je bodem \textit{\textbf{lok�ln�ho} maxima funkce $f$ na $M$}, pokud $\exists \delta > 0$ tak, �e $\forall X \in M \cap B(\overline{X}, \delta): f(\overline{X}) \geq f(X) $. Analogicky definujeme \textit{\textbf{lok�ln�} minimum funkce $f$ na mno�in� $M$}.
\end{e}

\begin{e}{Definice}{0}{Stacion�rn� bod}
Nech� $M \subseteq \mathbb{R}^n$ otev�en�, $f: M \rightarrow \mathbb{R}$, $\overline{X} \in M$. �ekneme, �e bod $\overline{X}$ je \textit{stacion�rn�m bodem funkce $f$}, pokud existuj� v�echny parci�ln� derivace funkce $f$ v bod� $\overline{X}$ a jsou nulov�.
\end{e}

\begin{e}{V�ta}{0}{Nutn� podm�nka existence lok�ln�ho extr�mu}
Pokud $a \in \mathbb{R}^n$ je bodem lok�ln�ho extr�mu funkce $F: \mathbb{R}^n \rightarrow \mathbb{R}$ a v $a$ existuj� v�echny parci�ln� derivace funkce $F$, potom jsou tyto nulov�.
\end{e}


\begin{e}{Definice}{0}{Klasifikace biline�rn�ch forem}
Nech� $F: \mathbb{R}^n \times \mathbb{R}^n \rightarrow \mathbb{R}$ je biline�rn� forma.
\begin{pitemize}
\item $F$ se naz�v� \textit{pozitivn� definitn�}, pokud $\exists \varepsilon > 0$ tak, �e $F(h, h) \geq \varepsilon\| h \|^2$, $\forall h \in \mathbb{R}^n$.
\item $F$ se \textit{naz�v� negativn� definitn�}, pokud je $-F$ pozitivn� definitn�.
\item $F$ se naz�v� \textit{indefinitn�}, pokud $F(g,g) < 0$ a $F(h,h) > 0$ pro n�jak� $g, h \in \mathbb{R}^n$.
\end{pitemize}
\end{e}

\begin{e}{Pozn�mka}{0}{0}
P�i ur�ov�n� toho, zda je biline�rn� forma pozitivn� definitn�, negativn� definitn�, nebo indefinitn� n�m m��e pomoci tzv. Sylvestrovo krit�rium, kter� tvrd� n�sleduj�c�:
\begin{pitemize}
\item jsou-li v�echny hlavn� subdeterminanty matice reprezentuj�c� biline�rn� formu $F$ kladn�, potom je $F$ pozitivn� definitn�.
\item jestli�e je prvn� hlavn� subdeterminant t�to matice z�porn� a pot� alteruj� znam�nka, je forma negativn� definitn�.
\item nenast�v�-li ani jedna z p�edchoz�ch dvou mo�nost� a v�echny hlavn� subdeterminanty jsou nenulov�, je $F$ indefinitn�.
\end{pitemize}
Pakli�e nenastane ��dn� z v��e uveden�ch mo�nost�, Sylvestrovo krit�rium n�m nepom��e a je nutno o typu biline�rn� formy rozhodovat jin�m zp�sobem (nap�. pomoc� vlastn�ch ��sel).
\end{e}

\begin{e}{V�ta}{0}{Posta�uj�c� podm�nka pro existenci lok�ln�ho extr�mu}
Nech� $G \subseteq \mathbb{R}^n$ je otev�en� mno�ina a $a \in G$. Nech� $F: G \rightarrow \mathbb{R}$ m� spojit� parci�ln� derivace druh�ho ��du. Jestli�e $Df(a) = 0$, potom plat�:
\begin{pitemize}
\item je-li $D^2f(a)$ pozitivn� definitn�, potom $a$ je bodem lok�ln�ho minima
\item je-li $D^2f(a)$ negativn� definitn�, potom $a$ je bodem lok�ln�ho maxima
\item je-li $D^2f(a)$ indefinitn�, potom v bod� $a$ nen� lok�ln� extr�m
\end{pitemize}
\end{e}

\begin{e}{V�ta}{0}{O v�zan�ch extr�mech (Lagrangeovy multiplik�tory)}
Nech� $G \subseteq \mathbb{R}^n$ je otev�en�. M�jme funkce $F, g_1, \ldots g_m,~m<n$, kter� maj� spojit� parci�ln� derivace. Zadefinujme mno�inu $M$ spole�n�ch nulov�ch bod� funkc� $g_i$, $i=1 \ldots m$, tedy:
$$M = \{ x \in \mathbb{R}^n: g_1(x) = \ldots = g_m(x) = 0\}$$
Je-li bod $a = [a_1,\ldots,a_n]$ bodem lok�ln�ho extr�mu funkce $F$ na $M$ a plat�-li, �e vektory $\nabla g_1(a), \ldots, \nabla g_m(a)$ jsou line�rn� nez�visl�, potom existuj� tzv. Lagrangeovy multiplik�tory $\lambda_1, \ldots, \lambda_m$ takov�, �e:
$$DF(a) + \lambda_1 Dg_1(a) + \ldots + \lambda_m Dg_m(a) = 0$$
neboli
$$\frac{\partial F}{\partial x_i}(a) = \sum_{k=1}^m\lambda_k\frac{\partial g_k}{\partial x_i}(a),~i=1,\ldots,n$$
\end{e}
\begin{e}{Report}{0}{Fiala} pozitivne definitne matice - trochu nechapem odkial beru tuto otazku, uz ju dali viac krat... totiz PDmatice su spomenute vramci otazky Parci�ln� derivace a tot�ln� diferenci�l, v�ty o st�edn� hodnot�, extr�my funkc� v�ce prom�nn�ch, v�ta o implicitn�ch funkc�ch. konkretne sa nimi urcuju extremy, a o moc viac som o nich nevedel :( nakonec za 2
\end{e}


\end{document}
